\originalchapter{作业七：不可逆过程与正则涨落}
\psource{8.044，2013 春，PS7，4 月 3 日截止。中文题面；据官方答案核对的 AI 独立解答。}
\originalsection{1．非理想气体自由膨胀}
\begin{exercise}
单原子经典气体 $P=N\kb T/(V-bN)-a(N/V)^2$，$a,b>0$，$C_V=3N\kb/2$ 不变。可忽略热容的绝热容器总容积 $V_0$，初在 $V_0/3$、温度 $T_i$；开孔使其不可逆扩至 $V_0$。(a) 哪些热力学量守恒；(b) 终态温度为何？
\end{exercise}
\begin{solution}
(a) $N,V_0$ 固定，外界热与功均零，气体内能守恒；$S$ 一般增大，$T,P$ 不守恒。(b) $U_V|_T=TP_T-P=aN^2/V^2$，故 $U=3N\kb T/2-aN^2/V+U_*$。等能条件给 $T_f=T_i-4aN/(3\kb V_0)$。$b$ 不显式进入能量式；若此结果进入相分离或非经典区，简单单相模型失效，不能接受负温终态。
\end{solution}
\originalsection{2．有限热源的 Carnot 取功}
\begin{exercise}
两个相同物体恒热容 $C_0$，初温 $T_H>T_C>0$，以无穷小可逆 Carnot 步骤运行至共同温度 $T_F$。(a) 求 $T_F$；(b) 求对外总功及其符号。
\end{exercise}
\begin{solution}
总熵不变：$C_0\ln(T_F/T_H)+C_0\ln(T_F/T_C)=0$，故 $T_F=\sqrt{T_HT_C}$。外界得功 $W_{\rm out}=C_0(T_H+T_C-2\sqrt{T_HT_C})=C_0(\sqrt{T_H}-\sqrt{T_C})^2>0$。
\end{solution}
\originalsection{3．盐冷却液氦}
\begin{exercise}
液氦初温 1.0 K，$C_L=aT^3$，$a=128\times10^4\,\mathrm{erg\,K^{-4}}$；盐初温 $T_0$，$C_S=bT^{-2}$，$b=15\times10^4\,\mathrm{erg\,K}$。系统机械、热隔离且无功。(a) 若终温 0.5 K，求 $T_0$；(b) 用 $a,b$ 保留形式求宇宙熵增。
\end{exercise}
\begin{solution}
能量守恒 $a[(0.5\,\mathrm K)^4-(1\,\mathrm K)^4]/4+b(1/T_0-1/(0.5\,\mathrm K))=0$，给 $T_0=0.25$ K。熵变
\[\Delta S=\frac a3[(0.5\,\mathrm K)^3-(1\,\mathrm K)^3]+\frac b2[T_0^{-2}-(0.5\,\mathrm K)^{-2}].\]
按数值以 K 为温度单位即 $-7a/24+6b$，为正；前项液体熵降，后项盐熵增，环境无变。
\end{solution}
\originalsection{4．扭摆测量的热噪声}
\begin{exercise}
扭摆势能 $K(\theta-\theta_0)^2/2$，转动动能 $I\dot\theta^2/2$，处于温度 $T$ 的经典平衡。(a) 一次测角的均方根不确定度；(b) 求 $\avg{\theta\dot\theta}$。
\end{exercise}
\begin{solution}
位置与角速度正则密度因子化；$\avg{(\theta-\theta_0)^2}=\kb T/K$，所以标准差 $\sqrt{\kb T/K}$。速度平均零且独立，$\avg{\theta\dot\theta}=0$。多次测量的误差还取决于相关时间，不能直接沿用单次值。
\end{solution}
\originalsection{5．氢原子的配分函数}
\begin{exercise}
束缚态 $\ket{n,l,m_l}$ 能量 $-A/n^2$，$A=13.6$ eV。(a) 求第一激发能级与基能级原子数比，计算 300 K、1000 K；注意态和能级的区别。(b) 即使不含非束缚态，证明配分函数发散；(c) 为什么统计力学在此似乎失效？
\end{exercise}
\begin{solution}
(a) 简并度 $\sum_{l=0}^{n-1}(2l+1)=n^2$；共同电子自旋因子相消，比值 $4e^{-3A/(4\kb T)}$，分别约 $1.78\times10^{-171}$、$1.57\times10^{-51}$（脚本使用 $\kb=8.617333262\times10^{-5}$ eV/K）。(b) $Z=\sum_{n\ge1}n^2e^{A/(n^2\kb T)}$，项甚至不趋零，必发散。(c) 无穷高 $n$ 的波函数占据无限体积；实际容器有限、原子相互作用与电离破坏孤立原子谱。应建立有限体积、包含自由电子和质子及适当平衡约束的模型，不能将未归一化“概率”当作统计力学失败。
\end{solution}
